Несмотря на то, что в рамках учебного проекта невозможно реализовать полноценный движок уровня Unreal Engine или Godot, нами была создана упрощённая, но функциональная архитектура, отвечающая базовым требованиям 2D"=геймдева и позволяющая эффективно реализовать задуманный игровой процесс.\cite{GameDev}

Движок построен на базе библиотеки SFML (Simple and Fast Multimedia Library) и написан на языке C++ с использованием принципов объектно"=ориентированного программирования. Основной акцент был сделан на модульность и переиспользуемость кода, что позволило разделять задачи между членами команды и упростить отладку.
На данный момент в состав движка входят следующие ключевые подсистемы:

Система рендеринга. Обеспечивает отрисовку спрайтов, фонов и пользовательского интерфейса. Поддерживается работа с текстурами, трансформациями (масштабирование, поворот, позиционирование) и базовыми шейдерами OpenGL. Все графические объекты наследуются от единого базового класса, что позволяет управлять ими единообразно.

Система пользовательского интерфейса (UI). Реализованы базовые виджеты "--- кнопки, текстовые метки, флажки, контейнеры. Все элементы UI автоматически обрабатывают события ввода (наведение, нажатие), что упрощает создание меню и игровых экранов. Поддерживается масштабирование под разные разрешения.

Система ввода. Инкапсулирует обработку событий клавиатуры и мыши. Поддерживается как polling-метод (опрос состояния), так и event-driven подход (реакция на события SFML). Это позволяет реализовывать как мгновенные действия (например, нажатие кнопки), так и длительные (перетаскивание провода).

Аудиосистема. Позволяет воспроизводить звуковые эффекты (SFX) и фоновую музыку. Используются .ogg"=файлы для баланса между качеством и размером. Реализованы простые настройки громкости и приглушения при переходе в меню.

Система локализации. Поддерживается два языка интерфейса "--- русский и английский. Все текстовые строки вынесены в JSON-файлы, что упрощает добавление новых языков в будущем.

Система управления настройками. Реализовано сохранение и загрузка пользовательских параметров: разрешение окна, ограничение FPS, уровень громкости. Данные хранятся в локальном конфигурационном файле.

Базовая физика и коллизии. Хотя игра не является физически симулируемой, для объектов проводов и устройств реализованы простые прямоугольные хитбоксы. Это позволяет определять пересечения при перетаскивании и проверять корректность соединений.

Управление памятью. Ввиду отсутствия сборщика мусора в C++, нами реализованы явные деструкторы и используются умные указатели (std::unique\_ptr) там, где это уместно. Однако централизованной системы управления ресурсами (например, Resource Manager) пока нет "--- все текстуры и звуки загружаются по мере необходимости и хранятся в памяти до завершения сцены, что является потенциальной точкой оптимизации.

Таким образом, движок представляет собой гибкий каркас, достаточный для выполнения поставленных задач проекта, и может быть расширен в будущем "--- например, за счёт кэширования ресурсов, сценичной архитектуры или системы частиц.

А поскольку в проекте формул не было, предлагаю просто 3 удивительные формулы, взятые мной из самых честных источников:

\begin{align}
IQ_{k+1} &=
\frac{IQ_k \cdot e^{-\lambda \cdot PPW}}{\Omega},\\[1em]
IQ_{\text{humanity}}(B) &=
IQ_0 - \int_{0}^{B} \text{tokens} \, \mathrm{d}B
- \text{SelfRespect},\label{eq:iq} \\[1em]
L(\theta) &=
\sum_{i=1}^{N}
\log\left(IQ_{\text{LLM}} - IQ_{\text{humanity}}\right)
+ \alpha \lVert TT \rVert_2 .\label{eq:l}
\end{align}

Например, формула ~\ref{eq:l} точно отражает влияние тик"=ток роликов на мозг человека. Нельзя не отметить, что уровень интеллекта человечества находится в постоянной зависимости от количества использованных токенов, что отражено в формуле ~\ref{eq:iq}.